% Part: sets-functions-relations
% Chapter: functions
% Section: basics

\documentclass[../../../include/open-logic-section]{subfiles}

\begin{document}

\olfileid{sfr}{fun}{bas}
\olsection{بنسټونه}

\begin{explain}
\emph{تابع} داسې نگاشت دے چې د يوه ورکړل شوي سټ هر !!{element}
د کوم ورکړل شوي (بل) سټ له يوه ټاکلي \psOblique{!!{element}} سره تړي۔
د بېلګې په توګه، د $1$ ورزياتولو عمل يوه تابع ټاکي: هر عدد~$n$
له يوه يوازيني عدد~$n+1$ سره تړل کېږي۔
  
په عمومي ډول، تابعې جوړې، درې ګونې او داسې نور د ننوتونو په توګه
اخيستلے شي او يو ډول وت بېرته ورکوي۔ له بنسټيز حسابه ډېرې تابعې
پېژنو۔ د بېلګې په توګه، جمع او ضرب تابعې دي۔ دوې دوه عددونه
اخلي او درېيم عدد بېرته ورکوي۔

په دې رياضيکي، تجريدي معنٰی، تابع يو \emph{تور بکس} دے:
يوازې دا خبره اهميت لري چې کوم وت له کوم ننوت سره تړل شوے دے،
نۀ دا چې وت په کومه طريقه حسابېږي۔
\end{explain}

\begin{defn}[تابع]
\emph{تابع} $f \colon A \to B$ د~$A$ د هر \psOblique{!!{element}}
د~$B$ له يوه \psOblique{!!{element}} سره تړون دے۔

$A$ د~$f$ \emph{د تعريف ساحه} او $B$ د~$f$ \emph{هدف سټ} بولو۔
د~$A$ !!{element}s د~$f$ ننوتونه يا \emph{ارګومېنټونه} بلل کېږي،
او د~$B$ هغه !!{element} چې~$f$ ئې له ارګومېنټ~$x$ سره تړي،
د ارګومېنټ~$x$ دپاره \emph{د~$f$ قيمت} بلل کېږي او په~$f(x)$ ليکل کېږي۔

د~$f$ \emph{د قيمتونو سټ} $\ran{f}$ د هدف سټ هغه فرعي سټ دے چې
د~$f$ له هغو قيمتونو جوړ دے چې د کوم ارګومېنټ دپاره ترلاسه کېږي؛
$\ran{f} = \Setabs{f(x)}{x \in A}$۔
\end{defn}

په \olref{fig:function} کښې انځوريز شکل د تابعو په پوهېدو کښې مرسته
کولے شي۔ کيڼې خواته بيضوي شکل د تابعې \emph{د تعريف ساحه} ښيي؛
ښي خوا ته بيضوي شکل د تابعې \emph{هدف سټ} ښيي؛ او غشے د تعريف
په ساحه کښې له يوه \emph{ارګومېنټ} څخه په هدف سټ کښې ورته
\emph{قيمت} ته ځي۔

\begin{figure}
  \olasset{assets/diagrams/function.tikz}
  \caption{تابع د يوه سټ د هر \psOblique{!!{element}} د بل سټ له
    يوه \psOblique{!!a{element}} سره تړون دے۔ غشے د تعريف په ساحه کښې
    له يوه ارګومېنټ څخه په هدف سټ کښې ورته قيمت ته ځي۔}
  \ollabel{fig:function}
\end{figure}

\begin{ex}
ضرب د طبيعي عددونو جوړې د ننوتونو په توګه اخلي او طبيعي عددونه
د وتونو په توګه ورکوي؛ نو له $\Nat \times \Nat$ (د تعريف له ساحې)
څخه $\Nat$ (هدف سټ) ته ځي۔ د قيمتونو سټ هم $\Nat$ دے، ځکه چې
هر $n \in \Nat$ د $n \times 1$ په توګه ترلاسه کېږي۔
\end{ex}

\begin{ex}
ضرب تابع ده، ځکه چې هر ننوت---د طبيعي عددونو هره جوړه---له يوه
يوازيني وت سره تړي: $\times \colon \Nat^2 \to \Nat$۔ خو له يوه
طبيعي عدد~$n$ سره داسې حقيقي عدد~$x$ تړل چې $x^2 = n$ وي،
تابع نۀ جوړوي، ځکه چې هر مثبت صحيح عدد~$n$ دوه مربع جذرونه لري:
$\sqrt{n}$ او~$-\sqrt{n}$۔ که يوازې غير منفي (اصلي) مربع جذر بېرته ورکړو،
تابع ترې جوړولے شو: $\sqrt{\phantom{X}} \colon \Nat \to \Real$۔
\end{ex}

\begin{ex}
هغه اړيکه چې د يوه ټولګي هر زده کوونکے د هغه له وروستۍ درجې سره
تړي، تابع ده---يو زده کوونکے په هماغه ټولګي کښې دوه بېلې وروستۍ
درجې نۀ شي اخيستلے۔ هغه اړيکه چې د يوه ټولګي هر زده کوونکے
د هغه له والدينو سره تړي، تابع نۀ ده: زده کوونکي کېدے شي هېڅ،
يا دوه، يا تر دوو زيات والدين ولري۔
\end{ex}

\begin{explain}
تابعې داسې تعريفولے شو چې په يوه کره طريقه د هر ممکن ارګومېنټ
دپاره د تابعې قيمت وټاکو۔ د دې بېلې لارې دا دي چې يو فورمول
ورکړو، د قيمت د حسابولو طريقه بيان کړو، يا د هر ارګومېنټ قيمت
په فهرست کښې وليکو۔ تابعې چې په هره طريقه تعريفوو، بايد ډاډمن
شو چې د هر ارګومېنټ دپاره يو، او يوازې يو، قيمت ټاکو۔
\end{explain}


\begin{ex}
فرض کړئ $f \colon \Nat \to \Nat$ داسې تعريف شوې ده چې $f(x) = x+1$
وي۔ دا تعريف~$f$ د داسې تابعې په توګه ټاکي چې طبيعي عددونه اخلي
او طبيعي عددونه ورکوي۔ تعريف وايي چې د ورکړل شوي طبيعي عدد~$x$
دپاره به~$f$ د هغه ځايناستے~$x+1$ ورکړي۔ په دې حالت کښې هدف
سټ $\Nat$ د~$f$ د قيمتونو سټ نۀ دے، ځکه چې طبيعي عدد~$0$ د
هېڅ طبيعي عدد ځايناستے نۀ دے۔ د~$f$ د قيمتونو سټ د ټولو مثبتو
صحيح عددونو سټ، $\Int^{+}$، دے۔
\end{ex}

\begin{ex}\ollabel{examplefunext}
فرض کړئ $g \colon \Nat \to \Nat$ داسې تعريف شوې ده چې $g(x) = x+2-1$
وي۔ دا وايي چې~$g$ يوه تابع ده چې طبيعي عددونه اخلي او طبيعي
عددونه ورکوي۔ د ورکړل شوي طبيعي عدد دپاره به~$g$ د~$x$ د
ځايناستي د ځايناستي مخکينے عدد، يعنې $x+1$، ورکړي۔
\end{ex}

\begin{explain}
اوس مو دوه تابعې، $f$ او $g$، له بېلو \emph{تعريفونو} سره په پام
کښې ونيولې۔ خو دا \emph{هماغه يوه تابع} ده۔ ځکه چې د هر طبيعي
عدد~$n$ دپاره $f(n) = n+1 = n+2-1 = g(n)$ لرو۔ په بله وينا:
د $f$ او~$g$ تعريفونه مو په بېلو معادلو هماغه يو نگاشت ټاکي۔ نو
په ضمني ډول د تابعو د قيمتونو له مخې د برابرۍ په يوه اصل تکيه کوو:
\[
  \text{که }\forall x\, f(x) = g(x)\text{، نو }f = g
\]
په دې شرط چې د $f$ او~$g$ د تعريف ساحه او هدف سټ يو شان وي۔
\end{explain}

\begin{ex}
تابعې د حالتونو په وېشلو هم تعريفولے شو۔ د بېلګې په توګه،
$h \colon \Nat \to \Nat$ داسې تعريفولے شو:
\[
h(x) =
\begin{cases}
  \frac{x}{2} & \text{که $x$ جفت وي} \\
  \frac{x+1}{2} & \text{که $x$ طاق وي۔}
\end{cases}
\]
ځکه چې هر طبيعي عدد يا جفت دے يا طاق، د دې تابعې وت به تل
طبيعي عدد وي۔ يوازې دا په ياد وساتئ چې که تابع د حالتونو په
وېشلو تعريفوئ، هر ممکن ننوت بايد په دقيقه توګه په يوه حالت کښې
راشي۔ ځينې وختونه د دې ثبوت ته اړتيا وي چې حالتونه ټول امکانات
رانغاړي او يو له بل سره ګډ نۀ دي۔
\end{ex}

\end{document}
